% Page 097

\seite{97}

\margmark{6. Juni      V}%
om 7.\ bis 15.\ Juni einschließlich dauern die Pfingst=\ob
ferien. Das Wetter ist günstig gut. Die Heuernte\ob
ragt mit Riesenschritten, sie verspricht einen vortreff=\ob
lichen Ertrag.

\margmark{24. Juni     M}%
it dem heutigen Tage beginnen die Sommerferien.\ob
Der Unterricht beginnt, da der Festtag in die Pfingst=\ob
ferien fiel, am Mittwoch, den 9.\ Juli.

\margmark{1. Juli      D}%
as Rheinland ist frei von der Besatzung. In ver=\ob
gangener Mitternacht fanden überall Befreiung=\ob
und Feiern statt. Glockenläuten, Böllerschießen,\ob
weil die sich Freundstellung besonders Besuchern Kundgebung\ob
die Freiheit an. Heute mittags von 12--13 Uhr gaben\ob
sämtliche Kirchenglocken des ganzen deutschen Vater=\ob
landes die Freude ob die vierte Grenz stund. Nur die\ob
deutschen an der deutschen Masse sind noch nicht frei.\ob
Hoffentlich dürfen auch sie bald in ihrem Mutter=\ob
lande zurückkehren!

\margmark{10. Juli     D}%
er Kriegerverein der Pfarrei hielt am 6.\ Juli in\ob
Olzheim ein 2.\ Divisionsfest ab, es war gleichzeitig ein\ob
großes Volksfest. Das Kreisfest findet am 20.\ob
Juli statt.

\margmark{18. Juli     D}%
er Reichstag wurde gestern aufgelöst. Die Neu=\ob
wahl ist am 14.\ September.

\margmark{28. Juli     H}%
eute sollte anläßlich seiner Reise ins befreite Ge=\ob
biet der Herr Reichspräsident Generalfeldmarschall\ob
von Hindenburg unsere Bezirk statt Feier besuchen. Fuhr\ob
ein in Koblenz gestern abend gegen 24 Uhr vorgefahren und\ob
wurde bekanntgegeben, die weitere anläßlich einer\ob
Hindenburgspeisung nicht Befreiungsfeier. Menschen\ob
die ihm leben lassen müßten, hat der Reichspräsident\ob
die weitere seinen abgesagt und eine Erfrierungsklat=\ob
sche abgebrochen. Heute war in Koblenz vor an einer\ob
Rheinfeier ein Rollstuhl teil und ließ dann nach\ob
Köln zurück. Er kommt später nach Trier.
\pend
